\section{The fixed-shift cofactor ladder}
\label{sec:tpc61-cofactor-ladder}

We now use the literal equation to bound how much of one row can be
captured by a specified missing cofactor.  The resulting weight is harmonic
in \(r\), rather than uniform over the cofactor set.

\subsection{Coprimality of a realized cofactor}

Write every source row in its inherited unique form
\begin{equation}
 m=d\ell,
 \qquad d\in[D,2D],\quad \ell\in[L,2L]\cap\mathbb P,
 \label{eq:tpc61-row-factorization}
\end{equation}
where \(d\) is squarefree, \((d,h_0)=1\), and \(L>|h_0|\) for large
\(X\).  On the common high carrier, the TPC--45 support gives
\((d,r)=1\) and squarefreeness of \(dr\)
\citep{WangTPC45,WangTPC59,WangTPC60}.

\begin{lemma}[Full row--cofactor coprimality]
\label{lem:tpc61-full-row-cofactor-coprimality}
For every actual high atom,
\begin{equation}
 \boxed{(m,r)=1.}
 \label{eq:tpc61-m-r-coprime}
\end{equation}
Moreover \((m,h_0)=1\).
\end{lemma}

\begin{proof}
It remains only to exclude \(\ell\mid r\).  If this held, then
\(\ell\mid rp=mj+h_0\), while \(\ell\mid mj\); hence
\(\ell\mid h_0\), contradicting \(\ell>|h_0|\).  Thus
\((\ell,r)=1\), and the inherited relation \((d,r)=1\) proves
\eqref{eq:tpc61-m-r-coprime}.  Finally \((d,h_0)=1\) and
\(\ell>|h_0|\) imply \((m,h_0)=1\).
\end{proof}

\subsection{Exact rowwise progression}

Let \(\mathcal J_X\) denote the fixed orbit slice.  It is contained in an
integer interval of diameter at most \(C_JJ\), where \(C_J\) depends only
on the declared compact normalized orbit interval.  Direct residue
splitting may delete points from this interval and therefore only improves
the upper bounds below.

\begin{proposition}[Exact cofactor ladder]
\label{prop:tpc61-exact-cofactor-ladder}
Fix a source row \(m\) and a positive integer \(r\) which is realized by
at least one high atom.  There are integers \(j_0,p_0\) such that every
integral solution of
\begin{equation}
 rp-mj=h_0
 \label{eq:tpc61-fixed-row-equation}
\end{equation}
has the unique form
\begin{equation}
 \boxed{
 j=j_0+rt,
 \qquad
 p=p_0+mt,
 \qquad t\in\mathbb Z.}
 \label{eq:tpc61-row-cofactor-ladder}
\end{equation}
In addition,
\begin{equation}
 p\equiv h_0r^{-1}\pmod m,
 \label{eq:tpc61-prime-residue-class}
\end{equation}
and this is a reduced residue class modulo \(m\).
\end{proposition}

\begin{proof}
By \cref{lem:tpc61-full-row-cofactor-coprimality}, multiplication by
\(m\) is invertible modulo \(r\).  Hence
\(mj\equiv-h_0\pmod r\) has one residue class
\(j\equiv j_0\pmod r\).  Substitution of \(j=j_0+rt\) into
\eqref{eq:tpc61-fixed-row-equation} gives
\(p=p_0+mt\), and every such pair is an integral solution.  Conversely,
subtracting two solutions gives \(r(p-p')=m(j-j')\); coprimality of
\(m,r\) recovers the same parameterization.  Reducing the original
equation modulo \(m\) proves \eqref{eq:tpc61-prime-residue-class}.
The residue is reduced because both \(r\) and \(h_0\) are coprime to
\(m\).
\end{proof}

For the next estimates, let
\begin{equation}
 N_m(r)
 :=\#\left\{j\in\mathcal J_X:
 \begin{array}{l}
 p=(mj+h_0)/r\text{ is prime},\ p>y,\\[-2pt]
 p=\Pplus(mj+h_0),\text{ and there is a }\gamma\\[-2pt]
 \text{for which the common high masks pass}
 \end{array}\right\}.
 \label{eq:tpc61-Nmr}
\end{equation}
All conditions after integrality only delete points from the ladder.

\begin{corollary}[Deterministic rowwise capacity]
\label{cor:tpc61-deterministic-row-capacity}
Uniformly in every realized \(m,r\),
\begin{equation}
 \boxed{
 N_m(r)\le 1+C_J\frac Jr.}
 \label{eq:tpc61-trivial-row-capacity}
\end{equation}
\end{corollary}

\begin{proof}
The admissible \(j\)'s lie in one residue class modulo \(r\) inside an
interval of diameter at most \(C_JJ\).
\end{proof}

\subsection{The available one-dimensional sieve improvement}

The prime coordinate in \eqref{eq:tpc61-row-cofactor-ladder} lies in one
reduced residue class modulo \(m\).  Brun--Titchmarsh therefore improves
\eqref{eq:tpc61-trivial-row-capacity} by a logarithm when the ladder is
long.  We include the short-ladder case explicitly so that no expression
with \(\log(J/r)\le0\) is used.

\begin{proposition}[Brun--Titchmarsh row bound]
\label{prop:tpc61-brun-titchmarsh-row-bound}
There is a constant depending only on the fixed orbit interval such that
\begin{equation}
 \boxed{
 N_m(r)
 \ll
 \frac m{\varphi(m)}
 \left(
  1+\frac{J/r}{\log(2+J/r)}
 \right).}
 \label{eq:tpc61-BT-row-bound}
\end{equation}
On the inherited row support,
\begin{equation}
 \frac m{\varphi(m)}=X^{o(1)}.
 \label{eq:tpc61-m-over-phi-subpower}
\end{equation}
Thus the prime condition saves at most a logarithm in this
coefficient-blind row census; it does not change its fixed-power exponent.
\end{proposition}

\begin{proof}
Put \(T=J/r\).  If \(T\) is bounded by a sufficiently large constant
depending on \(C_J\), \eqref{eq:tpc61-trivial-row-capacity} is \(O(1)\),
which is dominated by the right side of \eqref{eq:tpc61-BT-row-bound}
because \(m/\varphi(m)\ge1\).

Suppose now that \(T\) exceeds that constant.  The prime values
\(p=p_0+mt\) lie in an interval of length at most \(C_JmT+O(m)\) and in
the reduced class \eqref{eq:tpc61-prime-residue-class}.  The interval form
of the Brun--Titchmarsh theorem gives
\begin{equation}
 \#\{p\text{ on this ladder}\}
 \ll
 \frac{mT}{\varphi(m)\log T};
 \label{eq:tpc61-BT-long-case}
\end{equation}
enlarging the interval by a fixed factor changes only the implied constant.
The largest-prime and remaining physical conditions delete primes, so the
same upper bound holds for \(N_m(r)\).  Combining the long and short cases
and replacing \(\log T\) by the uniformly positive
\(\log(2+T)\) proves \eqref{eq:tpc61-BT-row-bound}; see, for example,
\citet{MontgomeryVaughan2007}.

Finally the standard maximal-order estimate
\(n/\varphi(n)\ll\log\log(3n)\), together with
\(m=X^{O(1)}\), proves \eqref{eq:tpc61-m-over-phi-subpower}.
\end{proof}

\begin{corollary}[Dyadic missing-cofactor census]
\label{cor:tpc61-dyadic-missing-census}
For a set \(E_m(R)\subseteq[R,2R]\cap\N\), the number of high orbit
times on row \(m\) whose cofactor lies in \(E_m(R)\) is at most
\begin{equation}
 \ll
 \frac m{\varphi(m)}|E_m(R)|
 \left(
  1+\frac{J/R}{\log(2+J/R)}
 \right).
 \label{eq:tpc61-dyadic-missing-count}
\end{equation}
In fixed-power notation, if
\begin{equation}
 J=X^{\theta_J+o(1)},\quad
 R=X^{\theta+o(1)},\quad
 |E_m(R)|=X^{e(\theta)+o(1)},
 \label{eq:tpc61-dyadic-count-exponents}
\end{equation}
then the exponent in \eqref{eq:tpc61-dyadic-missing-count} is at most
\begin{equation}
 e(\theta)+(\theta_J-\theta)_+.
 \label{eq:tpc61-dyadic-capacity-exponent}
\end{equation}
\end{corollary}

\begin{proof}
Sum \eqref{eq:tpc61-BT-row-bound} over \(r\in E_m(R)\), using
\(r\asymp R\).  Equations
\eqref{eq:tpc61-m-over-phi-subpower} and
\eqref{eq:tpc61-dyadic-count-exponents} give the exponent formula.
\end{proof}

The estimate is an upper capacity, not a lower bound for actual missing
mass.  Nevertheless it identifies the strongest conclusion obtainable
from the fixed congruence and a one-dimensional upper-bound sieve alone.
Any further fixed-power improvement must use the actual retained relation,
the physical weights, a lower exposure estimate, or cancellation.
